%% kapa: exam=1 round=30 year=2019 period=1 form=A subject=civil_law subject_ko=민법 seq=25 number=25
%% source: Q-Net 감정평가사 기출문제 (cst00309) artlSeq=5205205 file=extracted/1008/5205205/2019년 제30회 감정평가사 1차시험 문제지/2019년 제30회 감정평가사 1차시험 문제지/2019년도 제30회 감정평가사 1차 1교시 A형.hwp
%% conversion: kordoc 4.24.51 -> tools/gen1.py (auto); review=proofread-sonnet; flags=none
\begin{question}{25}
등기의 추정력에 관한 설명으로 옳은 것을 모두 고른 것은? (다툼이 있으면 판례에 따름)
\begin{framedtext}
ㄱ. 가등기가 그 등기명의인의 의사에 기하지 아니하고 위조된 서류에 의하여 부적법하게 말소된 사실이 인정되는 경우, 그 가등기는 여전히 적법하게 이루어진 것으로 추정된다. \par
ㄴ. 등기명의자가 허무인(虛無人)으로부터 소유권이전등기를 이어받았다는 사실만으로는 그 등기명의자가 적법한 권리자라는 추정은 깨트려지지 않는다. \par
ㄷ. 소유권이전등기의 원인으로 주장된 계약서가 진정하지 않은 것으로 증명된 경우에 그 등기의 적법추정은 복멸되는 것이고, 계속 다른 적법한 등기원인이 있을 것으로 추정할 수는 없다.
\end{framedtext}
\begin{choices}
  \item ㄱ
  \item ㄴ
  \item ㄷ
  \item ㄱ, ㄷ
  \item ㄱ, ㄴ, ㄷ
\end{choices}
\end{question}
\answer{4}
